\documentclass[11pt]{article}

\usepackage[margin=1in]{geometry}
\usepackage{amsmath,amssymb}
\usepackage{enumitem}
\usepackage[hidelinks]{hyperref}

\setlength{\parindent}{0pt}
\setlength{\parskip}{0.7em}

\title{My Current Rough Understanding of a Transformer Language Model}
\author{}
\date{}

\begin{document}

\maketitle

This is my current very rough understanding about the ``transformer'' thing.
I have yet to understand many other details. For now, I am focusing on the
high-level structure of this giant monster.

The note below is meant to record the basic architecture first. I will refine
it later after I learn more.

\section*{High-level architecture of a Transformer language model}

\begin{enumerate}[leftmargin=*, label=\textbf{\arabic*.}]

\item An input text is \textbf{tokenized} into a sequence of tokens:
\[
  t_1,\ t_2,\ \ldots,\ t_n.
\]
For now, a token can be thought of roughly as a word.

\item There is a learnable embedding matrix:
\[
  E \in \mathbb{R}^{N_{\mathrm{vocab}} \times D}.
\]
Each token corresponds to one row of \(E\). Therefore, every possible token is
associated with a vector in
\[
  \mathbb{R}^{D}.
\]
Initially, the entries of \(E\) are usually small random numbers. They become
useful through training.

\item Each input token \(t_i\) is mapped to a row vector:
\[
  x_i^T \in \mathbb{R}^{1 \times D}.
\]
Putting these row vectors together gives
\[
  X^{(0)} \in \mathbb{R}^{n \times D}.
\]
Positional information is also added so that the model knows the order of the
tokens.

\item From the current matrix \(X^{(0)}\), the model creates three matrices:
\[
  Q \qquad \text{(query)},
\]
\[
  K \qquad \text{(key)},
\]
\[
  V \qquad \text{(value)}.
\]
The query--key--value language comes from earlier attention used with RNN
encoder--decoder models.\footnote{In that older setting, an encoder represented
an input sequence, and a decoder generated an output sequence one token at a
time while using weighted combinations of the encoder representations. I will
treat that historical machinery as a black box here. The words ``query,''
``key,'' and ``value'' name computational roles; they are not literal database
objects.}
These are combined to form a score matrix \(S\). The score matrix is then
converted into an attention matrix \(A\).

The attention matrix describes how much each token should use information from
the other permitted tokens.

Using \(A\) and \(V\), the model produces a new matrix \(Y\). Each row of \(Y\)
still corresponds to one token, but now incorporates information from other
tokens in the text.

\item Attention is combined with a few additional neural-network operations to
produce
\[
  X^{(1)}.
\]
The matrix \(X^{(1)}\) has the same overall size as \(X^{(0)}\), but its token
representations contain more contextual information.

\item This block of operations is repeated several times:
\[
  X^{(0)}
  \longrightarrow
  X^{(1)}
  \longrightarrow
  X^{(2)}
\]
\[
  \longrightarrow
  \cdots
  \longrightarrow
  X^{(L)}.
\]
The final matrix \(X^{(L)}\) contains highly contextualized representations of
all the input tokens.

\item To predict the next token, the model takes the last row of \(X^{(L)}\)
and applies a learnable affine transformation.

This produces one score for every possible token in the vocabulary. The scores
are converted into probabilities, and the model chooses or samples the next
token.

\end{enumerate}

\section*{The whole architecture}

In a very compressed form, the whole architecture is:
\[
\begin{array}{c}
\text{input text}
\\[0.5em]
\downarrow
\\[0.5em]
\text{sequence of tokens}
\\[0.5em]
\downarrow
\\[0.5em]
\text{embedding vectors}
\\[0.5em]
\downarrow
\\[0.5em]
\text{repeated Transformer blocks}
\\[0.5em]
\downarrow
\\[0.5em]
\text{scores for all possible next tokens}
\\[0.5em]
\downarrow
\\[0.5em]
\text{next token}.
\end{array}
\]

\section*{A short reminder to myself}

At this stage, the most important picture is:

\[
  \text{tokens}
  \longrightarrow
  \text{vectors}
  \longrightarrow
  \text{contextualized vectors}
  \longrightarrow
  \text{next-token probabilities}.
\]

The repeated Transformer blocks are where the representations become more and
more contextual. I do not need to understand every internal detail yet in order
to keep this high-level structure in mind.

\end{document}
